% =====================================================================
%  Module Data Lakehouse -- TP 1 : Ingestion et traitement de données structurées
%  Compilation : pdflatex (2 passes)
% =====================================================================
\documentclass[11pt,a4paper]{article}

\usepackage[utf8]{inputenc}
\usepackage[T1]{fontenc}
\usepackage[french]{babel}
\usepackage{lmodern}
\usepackage[babel=true]{microtype}

\usepackage[a4paper,margin=2.3cm,headheight=14pt]{geometry}
\usepackage{fancyhdr}
\usepackage{titlesec}
\usepackage{needspace}
\usepackage{enumitem}
\usepackage{amssymb}

\usepackage[table]{xcolor}
\usepackage{booktabs}
\usepackage{tabularx}
\usepackage{array}
\usepackage{tikz}
\usetikzlibrary{positioning,arrows.meta,babel}
\usepackage{listings}

\usepackage[htt]{hyphenat}
\usepackage[hidelinks]{hyperref}
\usepackage{xurl}
\usepackage{caption}
\captionsetup{font=small,labelfont=bf,labelsep=period}

\setlength{\parskip}{0.5em}
\setlength{\parindent}{0pt}
\setlength{\emergencystretch}{3em}

% ---------- Couleurs ----------
\definecolor{lhblue}{HTML}{1F4E79}
\definecolor{lhteal}{HTML}{0F766E}
\definecolor{lhred}{HTML}{B91C1C}
\definecolor{lhviolet}{HTML}{6D28D9}
\definecolor{lhgray}{HTML}{F3F4F6}
\definecolor{lhbronze}{HTML}{A16207}
\definecolor{lhsilver}{HTML}{64748B}
\definecolor{lhgold}{HTML}{CA8A04}

% ---------- Titres, en-têtes ----------
\titleformat{\section}{\needspace{6\baselineskip}\Large\bfseries\color{lhblue}}{\thesection.}{0.6em}{}
\titleformat{\subsection}{\large\bfseries\color{lhblue}}{\thesubsection}{0.6em}{}
\pagestyle{fancy}
\fancyhf{}
\fancyhead[L]{\small\color{lhblue}Module Data Lakehouse}
\fancyhead[R]{\small\color{lhblue}TP 1 -- Données structurées}
\fancyfoot[C]{\small\thepage}

% ---------- Encadrés : texte simple, sans boîte ----------
% Un titre en gras introduit le paragraphe ; seul « Attention » est en rouge.
\newenvironment{resultat}{\par\medskip\noindent\textbf{Résultat attendu :}\ }{\par\medskip}
\newenvironment{astuce}{\par\medskip\noindent\textbf{Astuce :}\ }{\par\medskip}
\newenvironment{attention}{\par\medskip\noindent\textbf{\textcolor{lhred}{Attention :}}\ }{\par\medskip}

% ---------- Code ----------
\lstset{
  basicstyle=\ttfamily\small, columns=fullflexible, keepspaces=true,
  breaklines=true, showstringspaces=false, upquote=true,
  frame=single, framesep=5pt, xleftmargin=6pt, xrightmargin=6pt,
  backgroundcolor=\color{lhgray}, rulecolor=\color{black!15},
  commentstyle=\color{black!55}\itshape,
  keywordstyle=\color{lhblue}\bfseries, stringstyle=\color{black!70},
  literate=
   {é}{{\'e}}1 {è}{{\`e}}1 {ê}{{\^e}}1 {à}{{\`a}}1 {ç}{{\c c}}1
   {ù}{{\`u}}1 {î}{{\^i}}1 {ô}{{\^o}}1 {É}{{\'E}}1 {Ê}{{\^E}}1
}
\lstdefinestyle{shell}{language=bash,backgroundcolor=\color{black!85},
  basicstyle=\ttfamily\small\color{white},rulecolor=\color{black!85},
  commentstyle=\color{black!35}\itshape,keywordstyle=\color{white}}

\newcommand{\cmd}[1]{\texttt{#1}}
\newcommand{\cocher}{$\square$\ }
\pdfstringdefDisableCommands{\def\cmd#1{#1}}

\newenvironment{notion}{\par\medskip\noindent\textbf{\textcolor{lhblue}{Notion clé :}}\ }{\par\medskip}
\newcommand{\nb}[1]{\textcolor{lhblue}{\texttt{#1}}}

% =====================================================================
\begin{document}
% =====================================================================

\begin{tikzpicture}[remember picture,overlay]
  \fill[lhblue] (current page.north west) rectangle ([yshift=-5.2cm]current page.north east);
  \node[anchor=north west,text=white,font=\sffamily] at ([xshift=2.3cm,yshift=-1.3cm]current page.north west)
    {Module Data Lakehouse -- Cycle ingénieur 4\textsuperscript{e}/5\textsuperscript{e} année};
  \node[anchor=north west,text=white,font=\sffamily\bfseries\fontsize{26}{30}\selectfont] at ([xshift=2.3cm,yshift=-2.2cm]current page.north west)
    {TP 1 -- Données structurées};
  \node[anchor=north west,text=white!85!lhblue,font=\sffamily\large] at ([xshift=2.3cm,yshift=-3.6cm]current page.north west)
    {De la base de données au tableau de bord, en passant par Bronze, Silver et Gold};
\end{tikzpicture}
\thispagestyle{empty}
\vspace*{4.2cm}

{\large\bfseries\color{lhblue}En bref}\par\smallskip
\begin{tabularx}{\linewidth}{@{}>{\bfseries}l X@{}}
\toprule
Objectif & Construire un pipeline Lakehouse complet sur les données réelles d'un site
           de e-commerce, puis en tirer des indicateurs et un tableau de bord. \\
Durée    & 2 séances de 3 h (+ environ 2 h de travail personnel). \\
Prérequis & Support 0 terminé : \cmd{check\_env.py} affiche \og ENVIRONNEMENT PRÊT \fg{}. \\
À rendre & Un compte rendu, les notebooks exécutés et le tableau de bord
           (section~\ref{sec:livrables}). \\
\bottomrule
\end{tabularx}

% =====================================================================
\section{Le scénario}
% =====================================================================

Vous rejoignez l'équipe data d'\textbf{Olist}, une place de marché brésilienne qui permet à de petits
commerçants de vendre sur les grands sites de e-commerce. La direction vous pose quatre questions :

\begin{enumerate}[leftmargin=*,nosep]
  \item Comment évoluent nos ventes, mois par mois ?
  \item Quelles catégories de produits rapportent le plus ?
  \item Tenons-nous nos promesses de délai de livraison, et où avons-nous des problèmes ?
  \item Peut-on prévoir un retard de livraison dès la commande ?
\end{enumerate}

Les données sont réparties entre deux systèmes : l'application de vente, qui écrit dans une base
\textbf{PostgreSQL}, et des \textbf{fichiers CSV} exportés par d'autres services (clients, produits,
vendeurs). Votre mission : les réunir dans le Lakehouse, les rendre fiables, puis répondre aux questions.

% =====================================================================
\section{Les données}
% =====================================================================

Le jeu de données public \textbf{Olist} (Kaggle, licence CC BY-NC-SA 4.0) décrit environ
100\,000 commandes réelles passées entre 2016 et 2018. Il est fourni dans le dossier
\cmd{data/olist/} du dépôt.

\begin{center}
\small
\renewcommand{\arraystretch}{1.25}
\begin{tabularx}{\textwidth}{@{}>{\ttfamily}l l X r@{}}
\toprule
\normalfont\textbf{Table} & \textbf{Arrive par} & \textbf{Contenu} & \textbf{Lignes} \\
\midrule
orders          & PostgreSQL & Une ligne par commande : statut, dates (achat, livraison, estimation) & 99\,441 \\
order\_items    & PostgreSQL & Les articles de chaque commande : produit, vendeur, prix, frais de port & 112\,650 \\
order\_payments & PostgreSQL & Les paiements : moyen, nombre d'échéances, montant & 103\,886 \\
customers       & CSV & Les clients : identifiant, ville, état & 99\,441 \\
products        & CSV & Les produits : catégorie (en portugais), poids, dimensions & 32\,951 \\
sellers         & CSV & Les vendeurs : ville, état & 3\,095 \\
category\_translation & CSV & Traduction des catégories en anglais & 71 \\
\bottomrule
\end{tabularx}
\end{center}

\begin{attention}
Dans Olist, un \cmd{customer\_id} est créé \textbf{à chaque commande}. Pour identifier une même
personne d'une commande à l'autre, il faut utiliser \cmd{customer\_unique\_id}.
\end{attention}

% =====================================================================
\section{Ce que vous allez construire}
% =====================================================================

La figure~\ref{fig:pipeline} montre le pipeline complet : les tables de chaque couche, les
transformations entre les couches et le notebook qui réalise chaque étape. Il se lit en deux temps :
\textbf{de gauche à droite} sur la première ligne (ingestion), puis \textbf{de droite à gauche} sur la
seconde (préparation et usages).

\begin{figure}[ht]
\centering
\begin{tikzpicture}[
  font=\sffamily\small, >=Stealth,
  col/.style={draw=#1,thick,rounded corners=4pt,fill=#1!7,text width=3.75cm,align=left,
              inner sep=6pt,font=\sffamily\footnotesize},
  tag/.style={font=\sffamily\scriptsize\bfseries,text=white,fill=lhblue,rounded corners=2pt,inner sep=2.5pt},
  op/.style={font=\sffamily\scriptsize,text=black!70,align=center}]

% ---------------- Ligne 1 : ingestion ----------------
\node[col=gray] (pg) at (0,1.2) {\textbf{PostgreSQL}\\ schéma \texttt{ventes}\\[3pt]
  \texttt{orders}\\ \texttt{order\_items}\\ \texttt{order\_payments}};
\node[col=gray] (csv) at (0,-2.0) {\textbf{Fichiers CSV}\\ dossier \texttt{data/olist/}\\[3pt]
  \texttt{customers}\\ \texttt{products}\\ \texttt{sellers}\\ \texttt{category\_translation}};
\node[col=darkgray] (land) at (5.35,-2.0) {\textbf{Landing}\\ \texttt{s3://landing/}\\[3pt]
  un dossier par table\\ et par date de dépôt\\ fichiers non modifiés};
\node[col=lhbronze] (br) at (11.0,-0.4) {\textbf{Bronze} -- 7 tables\\[3pt]
  \texttt{orders}, \texttt{order\_items},\\ \texttt{order\_payments}\\[2pt]
  \texttt{customers}, \texttt{products},\\ \texttt{sellers},\\ \texttt{category\_translation}\\[4pt]
  \textit{colonnes ajoutées :}\\ \texttt{\_source}, \texttt{\_ingere\_le}};

\draw[->,thick] (pg.east) -- node[op,above]{JDBC (lot 1, puis lot 2)} (pg.east -| br.west);
\draw[->,thick] (csv) -- node[op,above]{dépôt} (land);
\draw[->,thick] (land.east) -- node[op,above]{lecture\\en texte} (land.east -| br.west);

% ---------------- Ligne 2 : préparation et usages ----------------
\node[col=lhsilver] (si) at (11.0,-7.1) {\textbf{Silver} -- tables propres\\[3pt]
  \texttt{customers}, \texttt{sellers}\\
  \texttt{products} \textit{+ catégorie EN}\\
  \texttt{orders} \textit{+ délais, retard}\\
  \texttt{order\_items}\\
  \texttt{order\_payments}};
\node[col=gray, below=1.0cm of si] (q) {\textbf{Quarantaine}\\[2pt]
  \texttt{quarantaine\_orders}\\ \texttt{quarantaine\_payments}};
\node[col=lhgold] (go) at (5.35,-7.1) {\textbf{Gold} -- prêtes à l'usage\\[3pt]
  \texttt{commandes}\\ \textit{\quad une ligne par commande}\\
  \texttt{ventes\_mensuelles}\\
  \texttt{ventes\_categorie}\\
  \texttt{livraison\_etat}\\
  \texttt{predictions\_retard}};
\node[col=lhblue,text width=3.4cm] (sql) at (0,-5.2) {\textbf{Requêtes SQL}\\ questions de la direction};
\node[col=lhblue,text width=3.4cm] (ml)  at (0,-7.1) {\textbf{Modèle ML}\\ scikit-learn : retards};
\node[col=lhblue,text width=3.4cm] (db)  at (0,-9.0) {\textbf{Tableau de bord}\\ Streamlit};

\draw[->,thick] (br.south) -- node[op,right,align=left]{typage, dédoublonnage,\\règles de qualité\\[2pt]\textit{lot 2 : \texttt{MERGE}}} (si.north);
\draw[->] (si.south) -- node[op,right]{lignes rejetées} (q.north);
\draw[->,thick] (si.west) -- node[op,above]{jointures,\\agrégats} (go.east);
\draw[->,thick] (go.west) -- ++(-0.5,0) |- (sql.east);
\draw[<->,thick] (go.west) -- (ml.east);
\draw[->,thick] (go.west) -- ++(-0.5,0) |- (db.east);

% ---------------- Notebooks ----------------
\node[tag,anchor=south east] at (pg.north east)  {01\_bronze};
\node[tag,anchor=south east] at (br.north east)  {01\_bronze};
\node[tag,anchor=south east] at (si.north east)  {02\_silver};
\node[tag,anchor=south east] at (go.north east)  {03\_gold\_sql\_incremental};
\node[tag,anchor=south east] at (ml.north east)  {04\_ml\_retards};
\node[tag,anchor=south east] at (db.north east)  {dashboard\_tp1.py};
\end{tikzpicture}
\caption{Pipeline du TP 1. Les étiquettes bleues indiquent le notebook qui réalise chaque étape.
Le modèle ML lit \texttt{gold.commandes} et y écrit \texttt{predictions\_retard}.}
\label{fig:pipeline}
\end{figure}

\begin{center}
\small
\renewcommand{\arraystretch}{1.3}
\begin{tabularx}{\textwidth}{@{}>{\bfseries}p{2.2cm} X X@{}}
\toprule
\normalfont\textbf{Couche} & \textbf{Rôle} & \textbf{Règle d'or} \\
\midrule
Bronze & Copie des sources, avec deux colonnes techniques : \cmd{\_source} et \cmd{\_ingere\_le} &
  On ne modifie rien. On ajoute, on n'écrase jamais l'historique. \\
Silver & Données nettoyées, typées, dédoublonnées et contrôlées &
  Une ligne invalide n'est pas supprimée en silence : elle va en quarantaine. \\
Gold & Tables prêtes à l'usage : une ligne par commande et des agrégats &
  Chaque indicateur se recalcule à partir de Silver. \\
\bottomrule
\end{tabularx}
\end{center}

Toutes les tables sont au format \textbf{Apache Iceberg} et enregistrées dans le catalogue sous
le nom \cmd{lakehouse.<couche>.<table>} : par exemple \cmd{lakehouse.silver.orders}.

% =====================================================================
\section{Déroulé}
% =====================================================================

\textbf{Avant de commencer, copiez les énoncés dans votre dossier personnel.} Dans JupyterLab, ouvrez un
terminal (\emph{File > New > Terminal}) et tapez :
\begin{lstlisting}[style=shell]
cp -r notebooks/tp1 apps/dashboard_tp1.py notebooks/mon_travail/
\end{lstlisting}
Travaillez \textbf{uniquement} dans \nb{notebooks/mon\_travail/}. Git ignore ce dossier : vos fichiers ne
seront jamais écrasés, et \cmd{git pull} récupérera sans conflit les mises à jour publiées par l'enseignant.
Ne faites ni \cmd{git commit} ni \cmd{git push} : le dépôt du cours est en lecture seule.

Les notebooks contiennent les consignes détaillées : les cellules marquées \cmd{\# À COMPLÉTER} sont à
écrire. Exécutez-les \textbf{dans l'ordre}.
Gardez sous la main l'\textbf{aide-mémoire PySpark} de la dernière page : il contient toutes les
opérations dont vous aurez besoin.

\begin{center}
\small
\renewcommand{\arraystretch}{1.3}
\begin{tabularx}{\textwidth}{@{}l >{\ttfamily}l X r@{}}
\toprule
& \normalfont\textbf{Notebook} & \textbf{Contenu} & \textbf{Durée} \\
\midrule
\textbf{Séance 1}
 & 01\_bronze & Chargement de la source, dépôt des fichiers, ingestion CSV et JDBC & 1 h \\
 & 02\_silver & Profilage, nettoyage, règles de qualité, quarantaine, intégrité & 2 h \\
\midrule
\textbf{Séance 2}
 & 03\_gold\_sql\_incremental & Tables Gold, questions SQL, lot 2 avec \cmd{MERGE}, \emph{time travel} & 1 h 30 \\
 & 04\_ml\_retards & Modèle de prédiction des retards, prédictions écrites en Gold & 1 h \\
 & apps/dashboard\_tp1.py & Tableau de bord Streamlit & 30 min \\
\bottomrule
\end{tabularx}
\end{center}

\begin{attention}
\textbf{Un seul notebook actif à la fois.} Chaque notebook ouvert garde sa propre session Spark,
qui réserve jusqu'à 2 Go de mémoire. Avant d'ouvrir le notebook suivant, arrêtez les précédents :
panneau \emph{Running Terminals and Kernels} (barre de gauche) > \emph{Shut Down All}.
Les tables restent dans le Lakehouse.
\end{attention}

\begin{astuce}
Avant chaque séance : \cmd{docker compose up -d}, puis \cmd{git pull} pour récupérer les dernières
versions. Si \cmd{git pull} refuse de s'exécuter, c'est qu'un fichier du cours a été modifié
par erreur : \cmd{git restore .} rétablit l'original, sans toucher à \nb{mon\_travail/}. Une erreur \og \texttt{Cellule à compléter : \dots} \fg{} signale simplement une partie que vous
n'avez pas encore écrite.
\end{astuce}

% =====================================================================
\section{Partie par partie}
% =====================================================================

\subsection*{Partie 1 -- Bronze \hfill \nb{01\_bronze}}

Vous chargez la base PostgreSQL dans l'état où elle était le \textbf{1\textsuperscript{er} juillet 2018}
(lot 1), vous déposez les fichiers CSV dans la zone \cmd{s3://landing/}, puis vous écrivez sept tables Bronze.

\begin{notion}
\textbf{On lit les fichiers bruts en texte.} Laisser Spark deviner les types (\cmd{inferSchema})
transforme le code postal \cmd{09790} en nombre \cmd{9790} : le zéro est perdu pour toujours.
Les types sont fixés plus tard, en Silver, de façon explicite.
\end{notion}

\begin{resultat}
\cmd{SHOW TABLES IN lakehouse.bronze} liste 7 tables. \cmd{bronze.orders} contient environ
86\,600 commandes : les autres n'existent pas encore au 1\textsuperscript{er} juillet 2018.
\end{resultat}

\subsection*{Partie 2 -- Silver \hfill \nb{02\_silver}}

Vous mesurez d'abord les défauts (valeurs nulles, doublons), puis vous écrivez une table Silver par
source. Un exemple complet est fourni pour \cmd{customers} ; vous traitez les autres sur le même modèle.

\begin{notion}
\textbf{Une valeur nulle n'est pas forcément une erreur.} Une commande en cours n'a pas de date de
livraison : c'est normal. Une commande \og livrée \fg{} sans date de livraison est une erreur :
elle part en \textbf{quarantaine}. Entre les deux, une incohérence sans gravité devient une
\textbf{alerte} : on la garde et on la signale.
\end{notion}

\begin{resultat}
Le bilan final compare Bronze et Silver table par table. Chaque écart s'explique par une règle :
doublons retirés, lignes en quarantaine.
\end{resultat}

\subsection*{Partie 3 -- Gold, SQL et ingestion incrémentale \hfill \nb{03\_gold\_sql\_incremental}}

Vous construisez \cmd{gold.commandes} (une ligne par commande) et trois tables d'indicateurs,
puis vous répondez en SQL aux questions de la direction. Vous chargez ensuite le \textbf{lot 2} :
deux mois et demi de nouvelles commandes et de changements de statut.

\begin{notion}
\textbf{Ingestion incrémentale.} On ne relit que les lignes modifiées depuis la dernière ingestion,
grâce à la colonne \cmd{updated\_at} (le \emph{filigrane}). L'instruction \cmd{MERGE INTO} applique
ces changements en Silver en une seule transaction : mise à jour des commandes existantes,
insertion des nouvelles. Iceberg garde l'ancienne version de la table : on peut la relire
avec \cmd{VERSION AS OF} (\emph{time travel}).
\end{notion}

\begin{resultat}
Après le \cmd{MERGE}, \cmd{silver.orders} contient environ 99\,400 commandes, et le nombre de
commandes \cmd{shipped} est divisé par deux environ : elles ont été livrées entre-temps.
\end{resultat}

\subsection*{Partie 4 -- Prédire les retards \hfill \nb{04\_ml\_retards}}

Vous entraînez deux modèles de classification (scikit-learn) qui estiment, au moment de la commande,
la probabilité d'un retard de livraison, puis vous comparez leurs performances. Les prédictions sont
écrites dans \cmd{gold.predictions\_retard} : le tableau de bord les affichera.

\begin{notion}
\textbf{Fuite de données.} Un modèle ne doit utiliser que des informations disponibles au moment de
la prédiction. Le notebook vous fait ajouter volontairement le délai réel de livraison :
le modèle devient presque parfait\dots{} et totalement inutilisable.
\end{notion}

\begin{resultat}
Une AUC ROC modeste, de l'ordre de 0,7 : les causes principales des retards (transporteur, météo,
grèves) ne figurent pas dans les données. Avec la fuite de données, l'AUC dépasse 0,95.
\end{resultat}

% =====================================================================
\section{Le tableau de bord}\label{sec:dashboard}
% =====================================================================

Ouvrez \nb{notebooks/mon\_travail/dashboard\_tp1.py} dans JupyterLab. Le chargement des données et le filtre de période
sont fournis. Complétez les cinq éléments demandés :

\begin{enumerate}[leftmargin=*,nosep]
  \item une ligne de quatre indicateurs : chiffre d'affaires, nombre de commandes, panier moyen, taux de retard ;
  \item la courbe du chiffre d'affaires mensuel ;
  \item le top 10 des catégories ;
  \item le taux de retard par état ;
  \item les commandes au plus fort risque de retard, selon le modèle de la partie 4, avec un seuil réglable.
\end{enumerate}

Pour le lancer, ouvrez un terminal dans JupyterLab (\emph{File > New > Terminal}) :
\begin{lstlisting}[style=shell]
streamlit run notebooks/mon_travail/dashboard_tp1.py --server.address 0.0.0.0
\end{lstlisting}
Puis ouvrez \url{http://localhost:8501} dans votre navigateur. N'utilisez pas l'adresse
\cmd{http://0.0.0.0:8501} affichée dans le terminal : elle n'est valable qu'à l'intérieur du conteneur.
Chaque enregistrement du fichier met la page à jour.

\begin{attention}
Le tableau de bord démarre sa propre session Spark. Sur une machine de 8 Go, arrêtez d'abord le
noyau des notebooks (\emph{Kernel > Shut Down Kernel}) pour libérer la mémoire.
\end{attention}

\begin{astuce}
Pensez à l'utilisateur : un titre clair, des unités (R\$, \%, jours), et des mois partiels exclus par
défaut. Un graphique qui fait croire à un effondrement des ventes en fin de période est une erreur
d'analyse, pas un détail.
\end{astuce}

% =====================================================================
\section{Livrables et évaluation}\label{sec:livrables}
% =====================================================================

À déposer sur l'espace de cours, dans une archive \cmd{NOM\_Prenom\_TP1.zip} :
\begin{itemize}[leftmargin=*,nosep]
  \item le \textbf{compte rendu} au format PDF (\cmd{NOM\_Prenom\_TP1.pdf}) ;
  \item les notebooks \cmd{01} à \cmd{04} \textbf{exécutés} ;
  \item le fichier \cmd{dashboard\_tp1.py}.
\end{itemize}

\subsection*{Le compte rendu}

Cinq à huit pages. Il ne recopie pas le code : il explique vos choix et vos résultats.
\begin{enumerate}[leftmargin=*,itemsep=2pt]
  \item \textbf{Architecture} : le schéma de votre pipeline et le rôle de chaque couche.
  \item \textbf{Qualité des données} : les défauts trouvés, les règles appliquées en Silver
        (bloquantes ou simples alertes) et le bilan Bronze / Silver.
  \item \textbf{Ingestion incrémentale} : le principe du filigrane et du \cmd{MERGE}, et ce que montre
        le \emph{time travel}.
  \item \textbf{Réponses à la direction} : les résultats des requêtes SQL, interprétés.
  \item \textbf{Modèle de prédiction} : les variables choisies, les métriques, la comparaison des deux
        modèles et la démonstration de la fuite de données.
  \item \textbf{Tableau de bord} : une capture d'écran et deux ou trois enseignements.
  \item \textbf{Réponses aux questions} des notebooks, rédigées et argumentées.
  \item \textbf{Limites et améliorations} : ce que vous feriez avec plus de temps, ou en production.
\end{enumerate}

\begin{center}
\small
\renewcommand{\arraystretch}{1.25}
\begin{tabularx}{0.9\textwidth}{@{}X r@{}}
\toprule
\textbf{Critère} & \textbf{Points} \\
\midrule
Bronze : ingestion CSV et JDBC correcte, colonnes techniques présentes & 3 \\
Silver : typage, dédoublonnage, règles de qualité, quarantaine & 4 \\
Gold : table des commandes et tables d'indicateurs correctes & 2 \\
Requêtes SQL et ingestion incrémentale (\cmd{MERGE}, \emph{time travel}) & 3 \\
Machine Learning : évaluation correcte, fuite de données expliquée, prédictions en Gold & 3 \\
Tableau de bord : complet, lisible, sans erreur d'analyse & 2 \\
Compte rendu : clarté, réponses argumentées, esprit critique & 3 \\
\midrule
\textbf{Total} & \textbf{20} \\
\bottomrule
\end{tabularx}
\end{center}

% =====================================================================
\section{En cas de problème}
% =====================================================================

\begin{center}
\small
\renewcommand{\arraystretch}{1.3}
\begin{tabularx}{\textwidth}{@{}>{\raggedright\arraybackslash}p{5.4cm}X@{}}
\toprule
\textbf{Symptôme} & \textbf{Solution} \\
\midrule
\texttt{Cellule à compléter : \dots} &
  Une partie de la cellule n'est pas encore écrite : remplacez l'appel \cmd{a\_completer(...)} par votre code. \\
\texttt{TABLE\_OR\_VIEW\_NOT\_FOUND} &
  La table vient d'une étape précédente non exécutée : relancez les cellules du notebook précédent, dans l'ordre. \\
\texttt{Fichier introuvable : .../data/olist/...} &
  Les données ne sont pas dans votre dossier : \cmd{git pull}, puis vérifiez \cmd{data/olist/}. \\
Le noyau meurt, ou \texttt{Java gateway process exited} &
  Mémoire insuffisante : \emph{Kernel > Restart}, fermez les autres notebooks et le tableau de bord. \\
\texttt{MERGE} : \texttt{cannot resolve ... columns} &
  Les colonnes de la source ne correspondent pas à la table cible : comparez avec \cmd{printSchema()}. \\
Vous avez exécuté le lot 2 deux fois &
  Sans danger : le chargement est idempotent. Pour tout recommencer, relancez le lot 1, puis les notebooks 01 et 02. \\
\url{http://0.0.0.0:8501} inaccessible &
  Ouvrez \url{http://localhost:8501}. \\
Le tableau de bord affiche une erreur Spark &
  Vérifiez que les tables Gold existent (\cmd{SHOW TABLES IN lakehouse.gold}) et que la session Spark des notebooks est arrêtée. \\
\bottomrule
\end{tabularx}
\end{center}

\clearpage
% =====================================================================
\section*{Annexe -- Aide-mémoire PySpark}\label{sec:aidememoire}
% =====================================================================

Imports utilisés partout : \cmd{from pyspark.sql import functions as F, Window}.

\begin{center}
\small
\renewcommand{\arraystretch}{1.22}
\begin{tabularx}{\textwidth}{@{}>{\raggedright\arraybackslash}p{4.1cm} >{\raggedright\arraybackslash\ttfamily\footnotesize}X@{}}
\toprule
\normalfont\textbf{Je veux\dots} & \normalfont\textbf{Code} \\
\midrule
\multicolumn{2}{@{}l}{\textbf{\textcolor{lhblue}{Lire}}} \\
Lire une table du Lakehouse & df = spark.table("lakehouse.silver.orders") \\
Lire des CSV (en texte) & spark.read.option("header", True).csv("s3a://landing/olist/sellers/") \\
Lire une table PostgreSQL & spark.read.format("jdbc").options(**JDBC).option("dbtable", "ventes.orders").load() \\
Regarder les données & df.show(5) \quad df.printSchema() \quad df.count() \\
\midrule
\multicolumn{2}{@{}l}{\textbf{\textcolor{lhblue}{Transformer}}} \\
Choisir et renommer des colonnes & df.select(F.col("a").alias("b"), "c") \\
Ajouter ou remplacer une colonne & df.withColumn("prix", F.col("prix").cast("double")) \\
Nettoyer un texte & F.trim("ville") \quad F.upper("etat") \quad F.initcap("ville") \quad F.lpad("cp", 5, "0") \\
Valeur conditionnelle & F.when(F.col("poids") > 0, F.col("poids")) \quad \textrm{(sinon : null)} \\
Remplacer une valeur nulle & F.coalesce(F.col("categorie"), F.lit("inconnue")) \\
Écart en jours entre deux dates & F.datediff("date\_fin", "date\_debut") \\
Filtrer des lignes & df.filter("statut = 'delivered'") \quad df.filter(F.col("prix") > 0) \\
Supprimer des colonnes & df.drop("\_source", "\_ingere\_le") \\
Supprimer les doublons & df.dropDuplicates(["order\_id"]) \\
Garder la ligne la plus récente & w = Window.partitionBy("id").orderBy(F.col("updated\_at").desc()) \newline
  df.withColumn("r", F.row\_number().over(w)).filter("r = 1").drop("r") \\
\midrule
\multicolumn{2}{@{}l}{\textbf{\textcolor{lhblue}{Combiner et agréger}}} \\
Jointure (gauche) & a.join(b, "order\_id", "left") \\
Lignes de A absentes de B & a.join(b, "product\_id", "left\_anti") \\
Lignes de A présentes dans B & a.join(b, "order\_id", "left\_semi") \\
Agréger par groupe & df.groupBy("order\_id").agg(F.sum("price").alias("total"), F.count("*").alias("nb")) \\
Valeur liée au maximum & F.max\_by("payment\_type", "payment\_value") \\
\midrule
\multicolumn{2}{@{}l}{\textbf{\textcolor{lhblue}{Écrire (tables Iceberg)}}} \\
Créer ou remplacer une table & df.writeTo("lakehouse.silver.x").using("iceberg").createOrReplace() \\
Ajouter des lignes & df.writeTo("lakehouse.bronze.x").append() \\
Utiliser SQL & spark.sql("SELECT ... FROM lakehouse.gold.x") \\
Exposer un DataFrame à SQL & df.createOrReplaceTempView("ma\_vue") \\
Historique d'une table & spark.sql("SELECT * FROM lakehouse.silver.x.snapshots") \\
\bottomrule
\end{tabularx}
\end{center}

\end{document}
